%! Exponential Functions — Unit 4, Lesson 4.3 — Exit Ticket (Answer Key)
\documentclass[10pt]{article}
\usepackage{precalculus-article}
\usepackage{precalculus-key}   % pulls in -boxes; do NOT also load -boxes
\newcommand{\TallMath}[1]{$\displaystyle #1\rule[-1.4em]{0pt}{3.2em}$}

\begin{document}

\pageheader{Unit 4, Lesson 4.3}{Exit Ticket --- Answer Key}
\namedateperiod

\vspace{0.16in}

\begin{enumerate}[itemsep=18pt]

\item For $f(x) = 3 \cdot (0.5)^x$, identify each of the following.

      \medskip
      Initial value $a$: \ans{$3$} \quad
      Base $b$: \ans{$0.5$} \quad
      Growth or decay? \ans{Decay}

\item State the end behavior of $g(x) = 5 \cdot 2^x$ using limit notation.

      \medskip
      $\displaystyle\lim_{x \to \infty}\, 5 \cdot 2^x = $ \ans{$\infty$}
      \qquad
      $\displaystyle\lim_{x \to -\infty}\, 5 \cdot 2^x = $ \ans{$0$}

\item A function has the following values.

      \renewcommand{\arraystretch}{1.4}
      \begin{tabular}{|c|c|c|c|c|}
        \hline
        $x$    & 0 & 1  & 2  & 3   \\
        \hline
        $f(x)$ & 6 & 18 & 54 & 162 \\
        \hline
      \end{tabular}

      \vspace{0.08in}
      Write $f(x)$ in the form $ab^x$. \quad $f(x) = $ \ans{$6 \cdot 3^x$}

      \smallskip
      {\small\color{charcoal}\textit{(Ratios: $18/6 = 3$, $54/18 = 3$, $162/54 = 3$; so $b = 3$ and $a = f(0) = 6$.)}}

\end{enumerate}

\end{document}
